\chapter{Em desenvolvimento}
\label{ApendiceC}
===== Em desenvolvimento
% A seguir são apresentados os parâmetros físicos do manipulador robótico, incluindo massas, dimensões e matrizes de inércia. Os valores encontram-se no Sistema Internacional (SI).

% \section{Massas e dimensões}
% As massas e dimensões físicas de cada componente do manipulador robótico foram medidas a partir do modelo tridimensional no \emph{Autodesk Fusion 360} e convertidas para o Sistema Internacional. 
% Esses parâmetros foram utilizados para a formulação da modelagem geométrica e para a determinação das matrizes de inércia, sendo posteriormente utilizados nas simulações computacionais. 
% A Tabela \eqref{tab:dimensoes_massas_m_kg} apresenta o comprimento, o raio, o diâmetro e a massa de cada componente, bem como o valor total da massa do manipulador robótico.

% \begin{table}[H]
% \centering
% \caption{Dimensões e massas dos componentes do manipulador (Sistema Internacional).}
% \label{tab:dimensoes_massas_m_kg}
% \begin{tabular}{l rrrr}
% \hline
% \textbf{Componente} & \textbf{Comprimento [m]} & \textbf{Raio [m]} & \textbf{Diâmetro [m]} & \textbf{Massa [kg]} \\
% \hline
% Base Geral   & 0{,}100 & 0{,}250  & 0{,}500 & 3{,}6530  \\
% Junta 1      & 0{,}005 & 0{,}0250 & 0{,}050 & 0{,}0265 \\
% Junta 2      & 0{,}005 & 0{,}0250 & 0{,}050 & 0{,}0265 \\
% Junta 3      & 0{,}005 & 0{,}0250 & 0{,}050 & 0{,}0265 \\
% Junta 4      & 0{,}005 & 0{,}0250 & 0{,}050 & 0{,}0265 \\
% Junta 5      & 0{,}005 & 0{,}0250 & 0{,}050 & 0{,}0265 \\
% Elo 1        & 0{,}100 & 0{,}0250 & 0{,}050 & 0{,}1115 \\
% Elo 2        & 0{,}300 & 0{,}0250 & 0{,}050 & 0{,}3129 \\
% Elo 3        & 0{,}300 & 0{,}0250 & 0{,}050 & 0{,}3129 \\
% Elo 4        & 0{,}300 & 0{,}0250 & 0{,}050 & 0{,}3129 \\
% Ferramenta   & 0{,}150 & 0{,}0075 & 0{,}015 & 0{,}0403 \\
% \hline
% \textbf{Total} & -- & -- & -- & \textbf{4{,}876} \\
% \hline
% \end{tabular}
% \end{table}


% \section{Momentos de inércia}
% Os momentos principais de inércia de cada componente do manipulador foram obtidos através do software \emph{Fusion 360} e estão apresentados na Tabela \eqref{tab:inercias_diag_kgm2}. 
% Os valores correspondem às componentes diagonais das matrizes de inércia ($I_{xx}$, $I_{yy}$ e $I_{zz}$), expressas em kg$\cdot$m$^{2}$. 
% Apesar das magnitudes reduzidas, esses parâmetros são relevantes para a análise dinâmica, principalmente em sistemas espaciais onde pequenas variações podem impactar o desempenho do controle.

% \begin{table}[H]
% \centering
% \caption{Momentos de inércia principais (valores diagonais) em kg$\cdot$m$^{2}$.}
% \label{tab:inercias_diag_kgm2}
% \begin{tabular}{l rrr}
% \hline
% \textbf{Componente} & $I_{xx}$ & $I_{yy}$ & $I_{zz}$ \\
% \hline
% Base Geral   & 0,07950   & 0,0795   & 0,144900 \\
% Junta 1      & 4,20$\times$10$^{-6}$ & 4,20$\times$10$^{-6}$ & 8,28$\times$10$^{-6}$ \\
% Junta 2      & 8,28$\times$10$^{-6}$ & 4,20$\times$10$^{-6}$ & 4,20$\times$10$^{-6}$ \\
% Junta 3      & 8,28$\times$10$^{-6}$ & 4,20$\times$10$^{-6}$ & 4,20$\times$10$^{-6}$ \\
% Junta 4      & 8,28$\times$10$^{-6}$ & 4,20$\times$10$^{-6}$ & 4,20$\times$10$^{-6}$ \\
% Junta 5      & 4,20$\times$10$^{-6}$ & 8,28$\times$10$^{-6}$ & 4,20$\times$10$^{-6}$ \\
% Elo 1        & 0,00137  & 0,00137  & 0,000597 \\
% Elo 2        & 0,02580   & 0,00174  & 0,025800 \\
% Elo 3        & 0,02580   & 0,00174  & 0,025800 \\
% Elo 4        & 0,02580   & 0,00174  & 0,025800 \\
% Ferramenta   & 7,82$\times$10$^{-5}$ & 1,62$\times$10$^{-6}$ & 7,82$\times$10$^{-5}$ \\
% \hline
% \end{tabular}
% \end{table}



% \section{Produtos de inércia}
% Além dos momentos principais, os produtos de inércia ($I_{xy}$, $I_{xz}$ e $I_{yz}$) também foram mantidos e encontram-se listados na Tabela
% \eqref{tab:inercias_prod_kgm2}. 
% Embora apresentem valores de ordem muito pequena, ao considerar mais casas decimais permite manter a consistência matemática dos modelos e evita simplificações que poderiam comprometer análises de precisão. 
% Essa escolha reforça a adequação do modelo a cenários de microgravidade, nos quais a dinâmica é sensível a perturbações de baixa intensidade.



% \begin{table}[H]
% \centering
% \caption{Produtos de inércia (valores fora da diagonal) em kg$\cdot$m$^{2}$.}
% \label{tab:inercias_prod_kgm2}
% \begin{tabular}{l rrr}
% \hline
% \textbf{Componente} & $I_{xy}$ & $I_{xz}$ & $I_{yz}$ \\
% \hline
% Base Geral   & 6,60$\times$10$^{-16}$ & 1,64$\times$10$^{-12}$ & -1,30$\times$10$^{-16}$ \\
% Junta 1      & -2,25$\times$10$^{-16}$ & 0 & 0 \\
% Junta 2      & -1,13$\times$10$^{-17}$ & 0 & 0 \\
% Junta 3      & -1,42$\times$10$^{-16}$ & 1,83$\times$10$^{-16}$ & -1,25$\times$10$^{-15}$ \\
% Junta 4      & 3,08$\times$10$^{-16}$ & -1,32$\times$10$^{-17}$ & -4,35$\times$10$^{-16}$ \\
% Junta 5      & -3,95$\times$10$^{-16}$ & 5,37$\times$10$^{-16}$ & -1,63$\times$10$^{-17}$ \\
% Elo 1        & 0 & -1,17$\times$10$^{-14}$ & 2,35$\times$10$^{-17}$ \\
% Elo 2        & 2,14$\times$10$^{-15}$ & -1,08$\times$10$^{-15}$ & 2,39$\times$10$^{-14}$ \\
% Elo 3        & 9,95$\times$10$^{-17}$ & -6,59$\times$10$^{-15}$ & 4,41$\times$10$^{-14}$ \\
% Elo 4        & 6,33$\times$10$^{-14}$ & 4,29$\times$10$^{-15}$ & -1,15$\times$10$^{-15}$ \\
% Ferramenta   & 5,43$\times$10$^{-15}$ & 4,39$\times$10$^{-16}$ & -5,66$\times$10$^{-16}$ \\
% \hline
% \end{tabular}
% \end{table}